\documentclass[11pt]{article}

\usepackage[margin=1in]{geometry}
\usepackage{amsmath,amssymb}
\usepackage{booktabs}
\usepackage{enumitem}
\usepackage{hyperref}
\hypersetup{colorlinks=true, linkcolor=blue, urlcolor=blue}

\title{Classification of Conditions and Problems\\ for an Equation of Motion}
\author{}
\date{}

\begin{document}
\maketitle

This note organizes the standard problem types associated with an
equation of motion (EOM) along two axes: (A) \emph{what data specifies
the problem} (the conditions imposed), and (B) \emph{what feature of
the solution's structure is being investigated}, once the EOM itself
is fixed.

\section*{A. By the Conditions Imposed}

\subsection*{1. Initial Value Problem (IVP / Cauchy problem)}
\begin{itemize}[leftmargin=*]
  \item \textbf{Condition:} the full state (position and velocity, or a
  phase-space point) at one instant $t_0$.
  \item \textbf{Governing results:} the Picard--Lindel\"of theorem gives
  local existence and uniqueness for Lipschitz forces; global existence
  requires a no-blow-up bound.
  \item This is what time-stepping integrators (RK4, symplectic
  schemes, etc.) actually solve.
\end{itemize}

\subsection*{2. Boundary Value Problem (BVP)}
\begin{itemize}[leftmargin=*]
  \item \textbf{Condition:} data at two or more distinct times/points,
  e.g. $x(t_0) = x_0$, $x(t_1) = x_1$ --- not a full state at one
  instant.
  \item Existence and uniqueness are \emph{not} automatic (unlike the
  IVP): zero, one, or many solutions can occur.
  \item \textbf{Prototypes:} the brachistochrone problem, geodesics
  between two points, Lambert's problem, and any Euler--Lagrange
  equation arising from a fixed-endpoint action principle.
  \item \textbf{Numerically:} shooting methods, relaxation /
  finite-difference methods, collocation.
\end{itemize}

\subsection*{3. Eigenvalue Problem (a BVP subtype)}
\begin{itemize}[leftmargin=*]
  \item A linear EOM together with boundary conditions admits
  nontrivial solutions only for special parameter values ---
  normal-mode frequencies, quantized energies.
\end{itemize}

\subsection*{4. Scattering Problem}
\begin{itemize}[leftmargin=*]
  \item \textbf{Condition:} an asymptotic free state as $t \to
  -\infty$ (or $r \to \infty$).
  \item \textbf{Sought:} the asymptotic outgoing state as $t \to
  +\infty$ --- the $S$-matrix or cross-section --- given a localized
  interaction region in between.
  \item Differs in kind from the IVP/BVP: the conditions sit at
  infinity, and often the trajectory itself is discarded --- only the
  in $\to$ out map matters.
\end{itemize}

\subsection*{5. Periodicity Condition}
\begin{itemize}[leftmargin=*]
  \item Search for $x(t+T) = x(t)$. Effectively a BVP with the two
  endpoints identified, but $T$ is itself unknown, making it a
  nonlinear eigenvalue-type problem.
  \item \textbf{Poincar\'e--Bendixson theorem} (2D flows): a bounded
  trajectory that does not approach a fixed point must approach a
  periodic orbit.
\end{itemize}

\subsection*{6. Inverse Problem}
\begin{itemize}[leftmargin=*]
  \item \textbf{Condition:} observed trajectories or scattering data.
  \item \textbf{Sought:} the force law / potential / EOM itself
  (inverse scattering; Newton inferring $1/r^2$ from Kepler's laws is
  the archetype).
\end{itemize}

\subsection*{7. Control / Optimization Problem}
\begin{itemize}[leftmargin=*]
  \item \textbf{Sought:} an external input, or a member of a family of
  EOMs, that steers the system to a target state while extremizing a
  cost functional --- Pontryagin's maximum principle.
  \item Structurally a BVP with an added adjoint (costate) system.
\end{itemize}

\section*{B. By the Question Asked About Solution Structure}

Given a fixed EOM, the following ask about the qualitative or global
structure of its solutions rather than about how the problem was
posed.

\begin{center}
\begin{tabular}{@{}p{4cm}p{5cm}p{6cm}@{}}
\toprule
\textbf{Problem} & \textbf{What it asks} & \textbf{Key tools / results} \\
\midrule
Existence / uniqueness / well-posedness &
Does a solution exist, is it unique, does it depend continuously on the data? &
Picard--Lindel\"of; Hadamard well-posedness \\
\addlinespace
Stationary solutions &
Fixed points $\dot{x} = f(x) = 0$, or static critical points $\nabla V = 0$ &
--- \\
\addlinespace
Linear stability &
Behavior under infinitesimal perturbation of a fixed point / orbit &
Jacobian eigenvalues (node / saddle / focus / center); Floquet multipliers for periodic orbits \\
\addlinespace
Nonlinear (Lyapunov) stability &
Does the full nonlinear flow stay near, or converge to, equilibrium? &
Lyapunov functions; LaSalle invariance principle \\
\addlinespace
Structural stability &
Does the qualitative phase portrait survive small changes to the EOM itself? &
Leads into bifurcation theory \\
\addlinespace
Bifurcation &
Qualitative change in fixed points / orbits as a parameter varies &
Saddle-node, pitchfork, transcritical, Hopf, period-doubling \\
\addlinespace
Attractors \& basins &
Long-time ($t \to \infty$) fate of trajectories in dissipative systems &
Fixed point / limit cycle / torus / strange attractor; basin of attraction \\
\addlinespace
Chaos &
Sensitive dependence on initial conditions &
Lyapunov exponents; Poincar\'e sections; homoclinic tangles (Melnikov method); routes to chaos (period-doubling, Ruelle--Takens, intermittency) \\
\addlinespace
Integrability &
Enough conserved quantities in involution to reduce the problem to quadratures &
Liouville integrability; action-angle variables; Noether symmetries; Lax pairs; Painlev\'e test \\
\addlinespace
Perturbation theory &
Fate of near-integrable systems under small perturbation &
Regular / singular perturbation; multiple scales; averaging; KAM theorem \\
\bottomrule
\end{tabular}
\end{center}

\section*{Structural Notes}

\begin{itemize}[leftmargin=*]
  \item \textbf{Conservative (Hamiltonian) systems have no attractors}
  --- Liouville's theorem forbids phase-space contraction --- so their
  long-time behavior is governed instead by the
  integrability/KAM/chaos axis (invariant tori surviving or breaking
  up into chaotic layers), not by attractor classification.
  \item \textbf{Dissipative systems} are where ``attractor'' language
  belongs; stability and bifurcation theory there describe how the
  attractor set itself changes.
  \item \textbf{Scattering is best read as an asymptotic BVP}: the
  conditions sit at $t = \pm\infty$ rather than at finite times, which
  is exactly why it needs its own machinery ($S$-matrix, asymptotic
  completeness) rather than shooting/relaxation methods.
  \item Bifurcation theory is the natural bridge between ``stability
  of a single solution'' and ``existence of qualitatively new families
  of solutions'' (e.g., a fixed point losing stability via a Hopf
  bifurcation and handing off to a limit cycle, which can itself
  period-double into chaos).
\end{itemize}

\end{document}
